% code_and_data.tex
\label{sec:code_and_data}

\subsection{Repository layout}

\begin{itemize}
\item \texttt{src/experiments/exp\_18\_tidal\_ionization.py}
  -- the load-bearing g-scan experiment.  Reproduced verbatim from
  Paper~I's
  \texttt{src/experiments/exp\_18\_tidal\_ionization.py}~\cite{menendez2026geometry};
  see the Provenance block in the script's docstring.  Sweeps
  gradient $g$ from $0$ (control) through $\sim\!5\times
  g_\text{crit}$, measures $T_\text{escape}(g)$ via the windowed
  $r_\text{pdf}$ diagnostic, and records the tidal-displacement
  signature ($x_e, x_p$ projections relative to system CoM).
\item \texttt{src/experiments/exp\_18\_tidal\_ionization.md} --
  Paper~III's expansion of Paper~I's stub: parameter table, expected
  runtime, PASS/PART/FAIL criteria for each audit-table row this
  experiment evidences.
\item \texttt{notes/plasma\_as\_gravitational\_ionization.md} --
  inherited from Paper~I with a Provenance prefix.  The
  Arnold-tongue / clock-density-gradient theoretical apparatus this
  paper builds on.
\item \texttt{notes/quantum\_roche\_limit.md} -- new in this paper.
  Synthesises Paper~I's gradient-ionization argument with the $M/d^3$
  Roche-scaling observation from Paper~I's follow-on catalogue
  entry~\#9, and outlines the $M_\text{min}(d)$ derivation that lands
  in the paper's main body.
\item \texttt{data/exp\_18\_tidal\_ionization.npy} -- g-scan output:
  rows of $(g, \text{escaped}, T_\text{escape}, x_e\_proj,
  x_p\_proj, r_\text{pdf,first\_window})$.  Written by the
  experiment script; tracked in git once the run produces it.
\item \texttt{data/exp\_18\_rpdf\_trajectories.npy} -- full
  $r_\text{pdf}(t)$ trajectories per $g$ value, for the trajectory
  figures.
\end{itemize}

\subsection{Engine}

The experiment script imports framework primitives
(\texttt{OctahedralLattice}, \texttt{CausalSession},
\texttt{enforce\_unity\_spinor}) from \texttt{dcl\_core.core} --
the continuous-amplitude submodule of the shared \texttt{dcl\_core}
package
(\href{https://github.com/JackDMenendez/dcl-core}{\nolinkurl{github.com/JackDMenendez/dcl-core}}),
which contains a verbatim port of Paper~I's
original \texttt{src/core/}.  The dependency is declared in
\texttt{virtual-env-requirements.txt} and installed automatically
by \texttt{setup.cmd} / \texttt{setup.sh}; no separate engine
provisioning is required.  At Paper~III's v1.0 release, the
\texttt{dcl\_core} dependency will be pinned to a specific
\texttt{@vX.Y.Z} tag (and that tag's Zenodo deposit DOI) so that
the released paper binds to an immutable engine version.
